\begin{table}[htbp]
\centering
\small
\caption{Шлюз переиспользования разборов: сверка на sentinel-выборке}
\label{tab:s13}
\begin{tabular}{>{\RaggedRight\hspace{0pt}\arraybackslash}p{0.460\linewidth}>{\RaggedRight\hspace{0pt}\arraybackslash}p{0.460\linewidth}}
\toprule
Этап разбора & Результат сверки на sentinel-выборке \\
\midrule
Stanza & совпал побитово на шести документах \\
NER & совпал побитово на шести документах \\
эмбеддинги & расхождение векторов до 3.5·10⁻⁷, признаков sem-v1 до 2.5·10⁻⁶ при допуске 1·10⁻⁴ \\
\bottomrule
\end{tabular}
\begin{minipage}{\linewidth}\footnotesize\vspace{2pt}
\textit{Примечание.} Единица анализа — этап разбора текста. Знаменатель: шесть документов sentinel-выборки, по три каждого класса. Данные: шлюз feature\_reuse\_gate, серия v2. Сверка построена на пересчёте с нуля и сравнении с кешем. Шкала: меньше расхождение значит надёжнее переиспользование. Сокращения: NER — распознавание именованных сущностей; sem-v1 — семейство семантических признаков. Пропуски: побитового совпадения у эмбеддингов не бывает: кеш считался пачками через границы документов, и порядок суммирования другой. Поправка на множественные сравнения не применялась.
\end{minipage}
\end{table}
